\chapter{Distributed Training Systems}\label{ch:distributed}

\Cref{ch:pretrain} ended with two uncomfortable numbers. A 7B-parameter model needs about
\SI{112}{GB} just for its weights, gradients and Adam moments---more than any single GPU holds---and
a frontier pretraining run needs $10^{24}$--$10^{26}$ FLOPs, which at a realistic
$4\times10^{14}$ FLOP/s per GPU is thousands of GPU-years. Both numbers force the same
conclusion: training must be spread across many devices, and the way it is spread decides
whether the cluster computes or waits.

Neither source book trains at this scale, so this chapter is new material. It is written for
someone who already knows what an all-reduce is and wants the \emph{cost model}: how each
parallelism strategy partitions memory, how many bytes it moves per step and over which links,
how those transfers overlap with compute, and how the strategies compose into the 3D/4D
layouts that real LLM training uses. Every algorithm is implemented from scratch in
\code{llmsys.dist} and tested in Lab~14 on CPU processes with the gloo backend, which exposes the
same collective API as NCCL.

\begin{figure}[t]
\centering
\begin{tikzpicture}[
  stage/.style={draw, rounded corners=2pt, minimum width=2.4cm, minimum height=0.95cm,
    align=center, font=\footnotesize, fill=#1},
  num/.style={circle, fill=sysblue, text=white, font=\scriptsize\bfseries, inner sep=1.2pt},
  arr/.style={-{Stealth[length=2mm]}, thick, sysblue!70}]
\node[stage=fcomm!15] (a) at (0,0) {Collectives \&\\ $\alpha$--$\beta$ model};
\node[stage=fdata!15] (b) at (2.9,0) {Data parallel\\ (bucketed DDP)};
\node[stage=fdata!15] (c) at (5.8,0) {ZeRO-1/2/3\\ \& FSDP};
\node[stage=fcompute!15] (d) at (8.7,0) {Tensor\\ parallel};
\node[stage=fcompute!15] (e) at (11.6,0) {Pipeline\\ parallel};
\node[stage=fmemory!20] (f) at (11.6,-1.8) {Sequence \&\\ context parallel};
\node[stage=fmemory!20] (g) at (8.7,-1.8) {Expert\\ parallel};
\node[stage=keygreen!12] (h) at (5.8,-1.8) {3D/4D layouts\\ \& planning};
\node[stage=warnred!12] (i) at (2.9,-1.8) {Reliability\\ at scale};
\foreach \x/\n in {a/1,b/2,c/3,d/4,e/5,f/6,g/7,h/8,i/9} \node[num] at (\x.north west) {\n};
\draw[arr] (a)--(b); \draw[arr] (b)--(c); \draw[arr] (c)--(d); \draw[arr] (d)--(e);
\draw[arr] (e)--(f); \draw[arr] (f)--(g); \draw[arr] (g)--(h); \draw[arr] (h)--(i);
\end{tikzpicture}
\caption{Roadmap. We first build the communication vocabulary (1), then partition the
\emph{batch} (2--3), the \emph{weights inside a layer} (4), the \emph{layers} (5), the
\emph{sequence} (6) and the \emph{experts} (7), and finally compose them (8) and keep the
resulting job alive (9).}
\label{fig:dist:roadmap}
\end{figure}

% ======================================================================
\section{The problem in numbers}\label{sec:dist:numbers}

Consider training a 70B dense model on 15T tokens, roughly the Llama~3 70B recipe. Training
compute is $6ND\approx6.3\times10^{24}$ FLOPs (\cref{eq:pretrain:6nd}). On H100s at 40\% MFU
($\approx\SI{396}{TFLOP/s}$ each) that is $1.6\times10^{10}$ GPU-seconds, about 500 GPU-years:
a 1-year run on 500 GPUs or a 2-month run on 3{,}000. Memory is just as binding: model states
alone are $16\times70\text{B}=\SI{1.12}{TB}$, fourteen 80\,GB GPUs before counting activations.

There are only a few things one can cut into pieces (\cref{fig:dist:axes}):
\begin{itemize}
\item the \textbf{batch}: every device holds the full model and processes different examples
  (data parallelism, DP); ZeRO/FSDP additionally shard the model states over the DP group;
\item the \textbf{weight matrices within a layer}: each device holds a slice of every matrix
  and computes a slice of every activation (tensor parallelism, TP);
\item the \textbf{stack of layers}: each device holds a contiguous group of layers and
  activations flow between devices (pipeline parallelism, PP);
\item the \textbf{sequence}: each device holds part of every sequence; attention needs the rest
  (sequence/context parallelism, SP/CP);
\item the \textbf{experts} of a mixture-of-experts layer (expert parallelism, EP).
\end{itemize}

\begin{figure}[t]
\centering
\begin{tikzpicture}[
  dev/.style={draw, fill=#1, minimum width=0.55cm, minimum height=0.55cm},
  title/.style={font=\footnotesize\bfseries, align=center},
  capt/.style={font=\scriptsize, align=center, text width=3.4cm}]
% --- DP
\begin{scope}[xshift=0cm]
\node[title] at (1.2,2.2) {Data parallel};
\foreach \g in {0,1,2,3} {
  \draw[fill=fcompute!25] (\g*0.65,0) rectangle ++(0.5,1.6);
  \draw[fill=fdata!50] (\g*0.65,-0.55) rectangle ++(0.5,0.4);
  \node[font=\tiny] at (\g*0.65+0.25,-0.35) {$x_\g$};
  \node[font=\tiny] at (\g*0.65+0.25,-0.85) {G\g};
}
\node[capt] at (1.2,-1.5) {full model on every GPU,\\ different batch slice;\\ all-reduce gradients};
\end{scope}
% --- TP
\begin{scope}[xshift=4.0cm]
\node[title] at (1.2,2.2) {Tensor parallel};
\foreach \g in {0,1,2,3} {
  \foreach \l in {0,1,2,3} \draw[fill=fcompute!25] (\g*0.65,\l*0.4) rectangle ++(0.5,0.35);
  \draw[fill=fdata!50] (\g*0.65,-0.55) rectangle ++(0.5,0.4);
  \node[font=\tiny] at (\g*0.65+0.25,-0.35) {$x$};
  \node[font=\tiny] at (\g*0.65+0.25,-0.85) {G\g};
}
\node[capt] at (1.2,-1.5) {every layer split column-\\ or row-wise; same batch;\\ all-reduce activations};
\end{scope}
% --- PP
\begin{scope}[xshift=8.0cm]
\node[title] at (1.2,2.2) {Pipeline parallel};
\foreach \g in {0,1,2,3} {
  \draw[fill=fcompute!25] (\g*0.65,\g*0.4) rectangle ++(0.5,0.35);
  \node[font=\tiny] at (\g*0.65+0.25,-0.85) {G\g};
}
\foreach \g in {0,1,2} \draw[-{Stealth[length=1.3mm]}, thick, fcomm] (\g*0.65+0.5,\g*0.4+0.18) -- (\g*0.65+0.65,\g*0.4+0.58);
\draw[fill=fdata!50] (0,-0.55) rectangle ++(0.5,0.4); \node[font=\tiny] at (0.25,-0.35) {$x$};
\node[capt] at (1.2,-1.5) {contiguous layer groups;\\ micro-batches stream\\ through; send/recv};
\end{scope}
% --- SP/CP
\begin{scope}[xshift=12.0cm]
\node[title] at (1.0,2.2) {Context parallel};
\foreach \g in {0,1,2,3} {
  \draw[fill=fcompute!25] (\g*0.55,0) rectangle ++(0.45,1.6);
  \draw[fill=fdata!50] (\g*0.55,-0.55) rectangle ++(0.45,0.4);
  \node[font=\tiny] at (\g*0.55+0.22,-0.35) {$x^{(\g)}$};
  \node[font=\tiny] at (\g*0.55+0.22,-0.85) {G\g};
}
\node[capt] at (1.0,-1.5) {sequence chunks;\\ exchange K/V by ring\\ or all-to-all};
\end{scope}
\end{tikzpicture}
\caption{What each parallelism cuts. Blue blocks are layers (a tall block is the full stack),
green boxes are the inputs each device sees. Expert parallelism (not shown) places different
experts of an MoE layer on different devices and routes tokens to them with all-to-all.}
\label{fig:dist:axes}
\end{figure}

They differ in what they save, what they cost in communication, and which network links that
communication can tolerate. To compare them we need a communication cost model.

% ======================================================================
\section{Communication fundamentals}\label{sec:dist:comm}

\subsection{Collectives}
All distributed training is built from a handful of collective operations over a group of
$p$ ranks (\cref{fig:dist:collectives}). Each rank starts with a buffer of $n$ bytes.

\begin{description}
\item[Broadcast] one rank's buffer is copied to all others.
\item[Reduce] the element-wise sum of all buffers ends up on one rank.
\item[All-reduce] every rank ends up with the element-wise sum. (Gradient averaging.)
\item[Reduce-scatter] the sum is computed, but rank $i$ only receives the $i$-th $1/p$ slice of it.
\item[All-gather] rank $i$ contributes slice $i$; every rank ends up with the concatenation.
\item[All-to-all] rank $i$ sends a distinct slice to every rank $j$ (a transpose of the
  data across ranks). Used by MoE dispatch and Ulysses sequence parallelism.
\end{description}
The identity \emph{all-reduce = reduce-scatter + all-gather} underlies both the ring
algorithm and ZeRO.

\begin{figure}[t]
\centering
\begin{tikzpicture}[scale=0.9, transform shape,
  c/.style={draw, minimum width=0.42cm, minimum height=0.42cm, font=\tiny, inner sep=0pt},
  lbl/.style={font=\scriptsize\bfseries}]
% all-reduce
\node[lbl] at (1.3,1.0) {All-reduce};
\foreach \r/\col in {0/fdata!30,1/fcompute!30,2/fmemory!40}
  \foreach \k in {0,1,2} \node[c, fill=\col] at (0+\k*0.42,0.3-\r*0.5) {$\r_\k$};
\draw[-{Stealth}, thick] (1.35,-0.2) -- (1.95,-0.2);
\foreach \r in {0,1,2} \foreach \k in {0,1,2} \node[c, fill=fcomm!25] at (2.2+\k*0.42,0.3-\r*0.5) {$\Sigma_\k$};
% reduce-scatter
\begin{scope}[xshift=4.6cm]
\node[lbl] at (1.3,1.0) {Reduce-scatter};
\foreach \r/\col in {0/fdata!30,1/fcompute!30,2/fmemory!40}
  \foreach \k in {0,1,2} \node[c, fill=\col] at (0+\k*0.42,0.3-\r*0.5) {$\r_\k$};
\draw[-{Stealth}, thick] (1.35,-0.2) -- (1.95,-0.2);
\foreach \r in {0,1,2} \node[c, fill=fcomm!25] at (2.2+\r*0.42,0.3-\r*0.5) {$\Sigma_\r$};
\end{scope}
% all-gather
\begin{scope}[xshift=9.2cm]
\node[lbl] at (1.3,1.0) {All-gather};
\foreach \r/\col in {0/fdata!30,1/fcompute!30,2/fmemory!40} \node[c, fill=\col] at (0+\r*0.42,0.3-\r*0.5) {$s_\r$};
\draw[-{Stealth}, thick] (1.35,-0.2) -- (1.95,-0.2);
\foreach \r in {0,1,2} {
  \node[c, fill=fdata!30] at (2.2,0.3-\r*0.5) {$s_0$};
  \node[c, fill=fcompute!30] at (2.62,0.3-\r*0.5) {$s_1$};
  \node[c, fill=fmemory!40] at (3.04,0.3-\r*0.5) {$s_2$};}
\end{scope}
% all-to-all
\begin{scope}[xshift=13.6cm]
\node[lbl] at (0.9,1.0) {All-to-all};
\foreach \r/\col in {0/fdata!30,1/fcompute!30,2/fmemory!40}
  \foreach \k in {0,1,2} \node[c, fill=\col] at (-0.6+\k*0.42,0.3-\r*0.5) {$\r_\k$};
\draw[-{Stealth}, thick] (0.6,-0.2) -- (0.95,-0.2);
\foreach \r in {0,1,2} {
  \node[c, fill=fdata!30] at (1.2,0.3-\r*0.5) {$0_\r$};
  \node[c, fill=fcompute!30] at (1.62,0.3-\r*0.5) {$1_\r$};
  \node[c, fill=fmemory!40] at (2.04,0.3-\r*0.5) {$2_\r$};}
\end{scope}
\end{tikzpicture}
\caption{Collective semantics for three ranks (rows). $r_k$ is chunk $k$ of rank $r$'s buffer,
$\Sigma_k=\sum_r r_k$. Reduce-scatter leaves rank $r$ with only $\Sigma_r$; all-to-all
transposes chunk ownership.}
\label{fig:dist:collectives}
\end{figure}

\subsection{The $\alpha$--$\beta$ cost model}
A point-to-point message of $n$ bytes costs
\begin{equation}
t(n) = \alpha + n\,\beta,
\end{equation}
where $\alpha$ is a fixed per-message latency (software, NIC, switch; a few to tens of
microseconds) and $\beta=1/B$ is the inverse bandwidth of the link. Small messages are
latency-bound, large ones bandwidth-bound; the crossover is at $n^\star=\alpha/\beta$, e.g.\
$\SI{10}{\micro s}\times\SI{50}{GB/s}=\SI{500}{KB}$ on an InfiniBand NDR link. Gradient
buckets of tens of MB are therefore firmly bandwidth-bound, which is why DDP bucketing
(\cref{sec:dist:dp}) is about \emph{overlap}, not latency.

\begin{proposition}[Bandwidth lower bound for all-reduce]
In an all-reduce over $p$ ranks, each rank must send at least $\frac{p-1}{p}n$ bytes and receive
at least $\frac{p-1}{p}n$ bytes for the reduce-scatter half, and the same again for the
all-gather half---$2\frac{p-1}{p}n$ in total per rank.
\end{proposition}
\emph{Sketch.} Each output element depends on $p$ inputs held on $p$ different ranks; for
rank $i$ to own its reduced $1/p$ slice, the other $p-1$ contributions to that slice
($\frac{p-1}{p}n$ bytes in total, after any partial pre-reduction) must reach it. Symmetric
arguments apply to distributing the result.

\subsection{Ring all-reduce}\label{sec:dist:ring}
The ring algorithm meets this bound. Arrange the ranks in a ring and split the buffer into $p$
chunks. In the \emph{reduce-scatter} phase, at step $s=0,\dots,p-2$ rank $r$ sends chunk
$(r-s)\bmod p$ to its right neighbor and receives chunk $(r-s-1)\bmod p$ from its left
neighbor, adding it to its own copy. After $p-1$ steps, rank $r$ holds the fully reduced chunk
$(r+1)\bmod p$. The \emph{all-gather} phase circulates the reduced chunks in the same pattern
without adding (\cref{fig:dist:ring}). Each of the $2(p-1)$ steps moves $n/p$ bytes on every
link simultaneously:
\begin{equation}
t_{\text{ring}} = 2(p-1)\,\alpha + 2\,\frac{p-1}{p}\,n\,\beta.
\label{eq:dist:ring}
\end{equation}
The bandwidth term is independent of $p$ for large $p$---the property that made ring
all-reduce (Patarasuk \& Yuan, 2009; popularized for deep learning by Baidu and Horovod)
the default for data parallelism. Its weakness is the latency term, linear in $p$; for small
messages or very large rings, tree or recursive halving/doubling algorithms with $O(\log p)$
latency win. NCCL chooses between ring, tree and (on NVSwitch systems) in-switch reduction
(NVLS) based on message size and topology.

\begin{figure}[t]
\centering
\begin{tikzpicture}[
  c/.style={draw, minimum width=0.5cm, minimum height=0.36cm, font=\tiny, inner sep=0pt},
  rk/.style={font=\scriptsize\bfseries}]
\def\SIG{S}
\foreach \step/\title/\data in {
  0/{start}/{{0,0,0,0},{1,1,1,1},{2,2,2,2},{3,3,3,3}},
  1/{after step 1}/{{0,0,0,30},{01,1,1,1},{2,12,2,2},{3,3,23,3}},
  2/{after step 2}/{{0,0,230,30},{01,1,1,301},{012,12,2,2},{3,123,23,3}},
  3/{after step 3 (done)}/{{0,S,230,30},{01,1,S,301},{012,12,2,S},{S,123,23,3}}} {
  \begin{scope}[xshift=\step*3.75cm]
    \node[font=\scriptsize] at (0.8,0.55) {\title};
    \foreach \row [count=\r from 0] in \data {
      \foreach \v [count=\k from 0] in \row {
        \edef\cur{\v}
        \ifx\cur\SIG \node[c, fill=fcomm!40] at (\k*0.52,-\r*0.42) {$\Sigma$};
        \else \node[c, fill=fdata!20] at (\k*0.52,-\r*0.42) {\v}; \fi
      }
      \node[rk, anchor=east] at (-0.3,-\r*0.42) {r\r};
    }
  \end{scope}
}
\end{tikzpicture}
\caption{Ring reduce-scatter on four ranks (rows) and four chunks (columns). Entries list
which ranks' contributions a chunk already contains; $\Sigma$ means fully reduced. After
$p-1=3$ steps rank $r$ owns the reduced chunk $(r+1)\bmod 4$. Three more steps of all-gather
copy each $\Sigma$ around the ring.}
\label{fig:dist:ring}
\end{figure}

\begin{lstlisting}[caption={Ring all-reduce from point-to-point messages (\code{llmsys/dist/collectives.py}, Lab 14 Task 1).},label={lst:dist:ring}]
def _sendrecv(send_buf, recv_buf, dst, src):
    reqs = [dist.isend(send_buf.contiguous(), dst), dist.irecv(recv_buf, src)]
    for r in reqs:                       # post both before waiting: no ring deadlock
        r.wait()

def ring_allreduce(t):
    rank, p = dist.get_rank(), dist.get_world_size()
    right, left = (rank + 1) % p, (rank - 1) % p
    chunks = list(torch.tensor_split(t.view(-1), p))     # views into t
    for s in range(p - 1):                               # reduce-scatter
        send_i, recv_i = (rank - s) % p, (rank - s - 1) % p
        buf = torch.empty_like(chunks[recv_i])
        _sendrecv(chunks[send_i], buf, right, left)
        chunks[recv_i] += buf
    owned = (rank + 1) % p
    for s in range(p - 1):                               # all-gather
        send_i, recv_i = (owned - s) % p, (owned - s - 1) % p
        buf = torch.empty_like(chunks[recv_i])
        _sendrecv(chunks[send_i], buf, right, left)
        chunks[recv_i].copy_(buf)
    return t
\end{lstlisting}

\begin{figure}[t]
\centering
\begin{tikzpicture}
\begin{loglogaxis}[width=0.85\linewidth, height=5.6cm, xmin=1e3, xmax=1e10,
  xlabel={message size $n$ (bytes)}, ylabel={time (s)}, label style={font=\small},
  tick label style={font=\scriptsize},
  legend style={font=\scriptsize, draw=none, at={(0.03,0.97)}, anchor=north west}]
% p=64, alpha=10us, beta=1/50GB/s
\addplot[very thick, warnred, domain=1e3:1e10, samples=60] {2*63*(1e-5 + x/50e9)};
\addplot[very thick, sysblue, domain=1e3:1e10, samples=60] {2*63*1e-5 + 2*63/64*x/50e9};
\addplot[very thick, keygreen, dashed, domain=1e3:1e10, samples=60] {2*6*(1e-5 + x/50e9)};
\addplot[very thick, napkin, densely dotted, domain=1e3:1e10, samples=60] {2*6*1e-5 + 2*63/64*x/50e9};
\legend{naive (via root), ring, binary tree, recursive halving/doubling}
\end{loglogaxis}
\end{tikzpicture}
\caption{$\alpha$--$\beta$ all-reduce cost for $p=64$, $\alpha=\SI{10}{\micro s}$,
$B=\SI{50}{GB/s}$ (analytic, \code{llmsys.dist.costs.allreduce\_time}). Tree algorithms win for
small messages (latency), ring and halving/doubling for large ones (bandwidth); the naive
root-based scheme is $\approx p/2$ times slower for large messages.}
\label{fig:dist:arcost}
\end{figure}

\subsection{The hardware the bytes travel on}
Inside an HGX H100 node, eight GPUs are fully connected through NVSwitch with NVLink~4: about
\SI{450}{GB/s} per GPU in each direction. Between nodes, each GPU typically has its own
\SI{400}{Gb/s} InfiniBand or RoCE NIC, i.e.\ \SI{50}{GB/s} per direction---an order of magnitude
less. PCIe Gen5 x16 sits at about \SI{64}{GB/s} per direction. The large intra-/inter-node
bandwidth gap is the single most important fact for parallelism layout: communication-heavy
strategies (TP, SP) stay inside the NVLink domain; communication-light ones (DP, PP) span
nodes. NCCL's reported \emph{bus bandwidth} for all-reduce rescales the measured algorithm
bandwidth by $2(p-1)/p$ (\cref{eq:dist:ring}) so that it can be compared directly with link
speed.

\begin{deepdive}[Inside NCCL]
NCCL splits a collective across several \emph{channels} (each a ring or tree mapped onto a
subset of links and driven by its own CUDA thread block) and pipelines each channel's chunk in
small slices so that sending, receiving and reducing overlap. It picks among three wire
\emph{protocols}: \emph{LL} (8-byte stores with an embedded flag: lowest latency, about half the
bandwidth), \emph{LL128} (128-byte lines with flags, near-full bandwidth on NVLink) and
\emph{Simple} (large chunks with explicit fences: full bandwidth, higher latency). On NVSwitch
systems, \emph{NVLS} performs the reduction inside the switch, and InfiniBand \emph{SHARP} does
the same in the network fabric, cutting the data each GPU must send for an all-reduce roughly in
half. Because NCCL kernels occupy SMs, heavy communication also steals a few percent of compute
from overlapping GEMMs---one reason comm/compute overlap is never perfectly free.
\end{deepdive}

\begin{napkinmath}[Gradient all-reduce for a 7B model]
bf16 gradients of 7B parameters are \SI{14}{GB}. Ring all-reduce among the 8 GPUs of one node:
$2\cdot\frac78\cdot\SI{14}{GB}/\SI{450}{GB/s}\approx\SI{54}{ms}$. Across 64 GPUs over IB at
\SI{50}{GB/s} per GPU: $\approx\SI{0.55}{s}$. Compare with the compute of one step: at 32k
tokens per GPU and 40\% MFU, $6\cdot7\text{B}\cdot32{,}768/\SI{396}{TFLOP/s}\approx\SI{3.5}{s}$.
Even across nodes, the all-reduce can be hidden behind the backward pass---if it is overlapped.
\end{napkinmath}

% ======================================================================
\section{Data parallelism}\label{sec:dist:dp}

Data parallelism replicates the model on $p$ ranks, gives each rank a different $1/p$ of the
global batch, and averages gradients before the optimizer step. If every rank starts from the
same weights and applies the same averaged gradient with the same optimizer, the replicas stay
identical forever, and the result is mathematically the same as single-device training on the
whole batch (for mean-reduced losses over equal-sized shards). Lab~14 checks this equivalence
to $10^{-6}$.

The simplest implementation all-reduces each \code{.grad} after \code{backward()}. It has two
problems. First, a model has hundreds of parameter tensors, many tiny (biases, norms), so the
step pays the latency $\alpha$ hundreds of times. Second, communication starts only after the
\emph{whole} backward pass, so the network idles during compute and the GPUs idle during
communication.

\paragraph{Bucketing and overlap.} Gradients become available in roughly reverse order of the
forward pass: the last layer's gradients are ready first. PyTorch DDP (Li et al., 2020) groups
parameters, in reverse order, into buckets of about \SI{25}{MB}, registers a hook on every
parameter, and launches an asynchronous all-reduce for a bucket as soon as all its gradients
exist. While NCCL reduces bucket $k$ on a separate CUDA stream, the backward pass keeps
computing the gradients of bucket $k+1$ (\cref{fig:dist:overlap}). Before the optimizer step,
the training loop waits for all outstanding reductions.

\begin{figure}[t]
\centering
\begin{tikzpicture}[x=0.75cm, y=0.6cm]
\node[font=\scriptsize\bfseries, anchor=east] at (-0.2,2.3) {no overlap};
\node[font=\scriptsize, anchor=east] at (-0.2,1.6) {compute};
\node[font=\scriptsize, anchor=east] at (-0.2,0.9) {network};
\draw[fill=fcompute!35] (0,1.3) rectangle (3,1.9) node[midway, font=\tiny] {forward};
\draw[fill=fcompute!60] (3,1.3) rectangle (9,1.9) node[midway, font=\tiny] {backward};
\draw[fill=fcomm!45] (9,0.6) rectangle (14,1.2) node[midway, font=\tiny] {all-reduce all gradients};
\draw[fill=fdata!45] (14,1.3) rectangle (15,1.9) node[midway, font=\tiny] {opt};
\begin{scope}[yshift=-2.6cm]
\node[font=\scriptsize\bfseries, anchor=east] at (-0.2,2.3) {bucketed};
\node[font=\scriptsize, anchor=east] at (-0.2,1.6) {compute};
\node[font=\scriptsize, anchor=east] at (-0.2,0.9) {network};
\draw[fill=fcompute!35] (0,1.3) rectangle (3,1.9) node[midway, font=\tiny] {forward};
\foreach \i/\lab in {0/B4,1/B3,2/B2,3/B1} \draw[fill=fcompute!60] (3+\i*1.5,1.3) rectangle ++(1.5,0.6) node[midway, font=\tiny] {bwd \lab};
\foreach \i/\lab in {0/b4,1/b3,2/b2,3/b1} \draw[fill=fcomm!45] (4.5+\i*1.5,0.6) rectangle ++(1.25,0.6) node[midway, font=\tiny] {AR \lab};
\draw[fill=fdata!45] (10.25,1.3) rectangle (11.25,1.9) node[midway, font=\tiny] {opt};
\draw[decorate, decoration={brace, amplitude=4pt, mirror}, warnred] (11.25,0.3) -- (15,0.3) node[midway, below=4pt, font=\scriptsize, warnred] {time saved};
\end{scope}
\end{tikzpicture}
\caption{Bucketing overlaps gradient communication with the remaining backward computation.
Only the reduction of the last bucket (the first layers) stays exposed. Buckets are formed in
reverse layer order because that is the order in which gradients appear.}
\label{fig:dist:overlap}
\end{figure}

\begin{lstlisting}[caption={Bucketed DDP from scratch (\code{llmsys/dist/ddp.py}, abridged).},label={lst:dist:ddp}]
class BucketedDDP(nn.Module):
    def __init__(self, module, bucket_cap_mb=25.0):
        super().__init__()
        self.module, self.world = module, dist.get_world_size()
        broadcast_parameters(module)                    # identical start everywhere
        self.buckets, cur, size = [], [], 0
        for p in reversed([p for p in module.parameters() if p.requires_grad]):
            cur.append(p); size += p.numel() * p.element_size()
            if size >= bucket_cap_mb * 2**20:
                self.buckets.append(_Bucket(cur)); cur, size = [], 0
        if cur: self.buckets.append(_Bucket(cur))
        for b in self.buckets:
            for p in b.params:
                self._bucket_of[p] = b
                p.register_post_accumulate_grad_hook(self._on_grad_ready)

    def _on_grad_ready(self, p):                         # runs inside backward()
        b = self._bucket_of[p]; b.ready += 1
        if b.ready == len(b.params):                     # bucket complete
            b.buffer = torch.cat([q.grad.reshape(-1) for q in b.params])
            b.work = dist.all_reduce(b.buffer, async_op=True)

    def finish_gradient_synchronization(self):
        for b in self.buckets:
            b.work.wait(); b.buffer /= self.world
            # ... copy slices of b.buffer back into each q.grad
\end{lstlisting}

\begin{systemslens}[What production DDP adds]
\begin{itemize}
\item \textbf{Gradient-as-bucket-view:} \code{.grad} tensors are views into the flat bucket
  buffer, so no \code{torch.cat}/copy-back is needed (our version copies twice).
\item \textbf{\code{no\_sync()}} for gradient accumulation: skip the all-reduce for all but the
  last micro-batch, otherwise accumulation multiplies communication by $k$.
\item \textbf{Unused parameters} (e.g.\ inactive MoE experts) never fire their hook; DDP must
  detect them (\code{find\_unused\_parameters}) or the bucket never completes---exactly the
  error our implementation raises.
\item \textbf{Bucket order} is re-learned after the first iteration from the actual order in
  which hooks fired.
\item \textbf{Compression hooks} (fp16/bf16 casting, PowerSGD) reduce bytes at some accuracy
  risk.
\end{itemize}
\end{systemslens}

\paragraph{How far can DP scale?} DP alone scales the batch: with $p$ replicas and a per-GPU
micro-batch of $b$ sequences, the global batch is $pb$ sequences. Beyond a \emph{critical batch
size} (McCandlish et al., 2018), larger batches stop reducing the number of steps needed, so
extra replicas waste compute. LLM pretraining typically uses global batches of 1--16M tokens.
And DP does nothing for memory: every replica still holds all 16 bytes per parameter. The next
section removes that redundancy.

\subsection{A step-time model for data parallelism}\label{sec:dist:dpmodel}
Let $t_c$ be the compute time of one step on one GPU and $t_m$ the time of the gradient
all-reduce. Without overlap, the step takes $t_c+t_m$; with perfect overlap,
$\max(t_c,t_m)$ (in practice the last bucket is always exposed). The \emph{scaling efficiency}
relative to a single GPU is therefore between
\begin{equation}
\eta_{\text{no overlap}}=\frac{t_c}{t_c+t_m}\qquad\text{and}\qquad
\eta_{\text{overlap}}=\frac{t_c}{\max(t_c,t_m)}.
\end{equation}
For ring all-reduce $t_m\approx 2\Psi/B$ is almost independent of $p$, while $t_c$ is
proportional to the \emph{local} batch. Weak scaling (fixed local batch, growing global batch)
therefore scales almost perfectly; strong scaling (fixed global batch split over more GPUs)
shrinks $t_c$ until communication dominates. \Cref{fig:dist:dpeff} plots both bounds for a 7B
model as the local batch shrinks.

\begin{figure}[t]
\centering
\begin{tikzpicture}
\begin{semilogxaxis}[width=0.85\linewidth, height=5.4cm, xmin=1024, xmax=131072, ymin=0, ymax=1.05,
  xlabel={tokens per GPU per step}, ylabel={DP scaling efficiency}, label style={font=\small},
  tick label style={font=\scriptsize}, log basis x=2,
  legend style={font=\scriptsize, draw=none, at={(0.98,0.03)}, anchor=south east}]
\addplot[very thick, sysblue, domain=1024:131072, samples=80] {(x*1.0617e-4)/((x*1.0617e-4)+0.56)};
\addplot[very thick, sysblue, dashed, domain=1024:131072, samples=80] {min(1,(x*1.0617e-4)/0.56)};
\addplot[very thick, warnred, domain=1024:131072, samples=80] {(x*1.0617e-4)/((x*1.0617e-4)+2.24)};
\addplot[very thick, warnred, dashed, domain=1024:131072, samples=80] {min(1,(x*1.0617e-4)/2.24)};
\legend{50 GB/s no overlap, 50 GB/s full overlap, 12.5 GB/s no overlap, 12.5 GB/s full overlap}
\end{semilogxaxis}
\end{tikzpicture}
\caption{Analytic DP efficiency for a 7B model (bf16 gradients, 14\,GB; ring all-reduce with
large $p$; compute at 40\% MFU of an H100, so $t_c\approx\SI{0.106}{ms}$ per token). With
overlap, efficiency stays at 1 until the local batch drops below $t_m/t_c$ tokens (about 5k
tokens at \SI{50}{GB/s}); without overlap it degrades smoothly. A slower network (e.g.\ 100\,Gb
Ethernet, \SI{12.5}{GB/s}) moves the knee $4\times$ to the right.}
\label{fig:dist:dpeff}
\end{figure}

\paragraph{Feeding data to many ranks.} Every DP rank must see a \emph{different} slice of
each global batch, and the union must be the same as in single-device training. PyTorch's
\code{DistributedSampler} does this by permuting indices with a shared seed per epoch
(\code{sampler.set\_epoch(e)}) and giving rank $r$ every $p$-th index starting at $r$. At
pretraining scale, data is pre-tokenized into many binary shards; each rank streams its own
subset of shards through a memory-mapped reader (\code{llmsys.data.TokenMemmapDataset}) and
records its position so that a resumed job continues exactly where it stopped. Forgetting
\code{set\_epoch} silently repeats the same order every epoch; giving two ranks the same seed
for shard selection silently duplicates data.

\paragraph{Global batch bookkeeping.} With micro-batch $b$ sequences of length $s$, gradient
accumulation $k$ and DP degree $d$, one optimizer step consumes $b\cdot k\cdot d\cdot s$ tokens.
Changing the number of GPUs while keeping the optimization identical means adjusting $k$ so the
product stays fixed; the learning rate schedule is defined over optimizer steps, not GPU-hours.

% ======================================================================
\section{ZeRO and fully sharded data parallelism}\label{sec:dist:zero}

In DP every rank stores identical copies of weights, gradients and optimizer moments. ZeRO
(Rajbhandari et al., 2020) observes that each rank only \emph{needs} the full set of some of
these, some of the time, and shards the rest over the $N_d$ data-parallel ranks in three
stages. With $\Psi$ parameters, 2-byte weights and gradients and $K=12$ bytes of fp32
optimizer state per parameter:
\begin{align}
\text{DP:}\quad & (2+2+K)\,\Psi \\
\text{ZeRO-1 (optimizer states):}\quad & (2+2+K/N_d)\,\Psi \\
\text{ZeRO-2 (+ gradients):}\quad & (2+(2+K)/N_d)\,\Psi \\
\text{ZeRO-3 / FSDP (+ parameters):}\quad & (2+2+K)\,\Psi/N_d
\end{align}

\begin{table}[t]
\centering\small
\begin{tabular}{lrrl}
\toprule
stage & bytes/param & 7.5B on 64 GPUs & communication per step \\
\midrule
DP (stage 0) & 16 & 120.0\,GB & all-reduce grads: $2\Psi$ \\
ZeRO-1 & $4 + 12/N_d$ & 31.4\,GB & reduce-scatter grads + all-gather params: $2\Psi$ \\
ZeRO-2 & $2 + 14/N_d$ & 16.6\,GB & same: $2\Psi$ \\
ZeRO-3 / FSDP & $16/N_d$ & 1.9\,GB & all-gather $\times2$ + reduce-scatter: $3\Psi$ \\
\bottomrule
\end{tabular}
\caption{Model-state memory per GPU (computed with \code{llmsys.dist.costs.zero\_memory\_per\_gpu})
and communication volume per rank (in units of the gradient size $\Psi$ bytes) for the ZeRO
stages.}
\label{tab:dist:zero}
\end{table}

\begin{figure}[t]
\centering
\begin{tikzpicture}[x=0.4cm, y=0.42cm]
\foreach \stg/\name/\w/\g/\o in {0/DP/1/1/1, 1/ZeRO-1/1/1/0, 2/ZeRO-2/1/0/0, 3/ZeRO-3/0/0/0} {
  \pgfmathsetmacro{\X}{\stg*9}
  \node[font=\footnotesize\bfseries] at (\X+3.6,7.4) {\name};
  \foreach \gp in {0,1,2,3} {
    \pgfmathsetmacro{\x}{\X+\gp*1.9}
    % weights: height 1 (2 bytes), grads: 1, optimizer: 6 (12 bytes)
    \ifnum\w=1 \draw[fill=fcompute!55] (\x,0) rectangle ++(1.5,1);
    \else \draw[fill=fcompute!55] (\x,0) rectangle ++(1.5,0.25); \fi
    \ifnum\g=1 \draw[fill=fcomm!45] (\x,1.2) rectangle ++(1.5,1);
    \else \draw[fill=fcomm!45] (\x,1.2) rectangle ++(1.5,0.25); \fi
    \ifnum\o=1 \draw[fill=fmemory!60] (\x,2.4) rectangle ++(1.5,6*0.75);
    \else \draw[fill=fmemory!60] (\x,2.4) rectangle ++(1.5,1.5*0.75); \fi
    \node[font=\tiny] at (\x+0.75,-0.5) {G\gp};
  }
}
\draw[fill=fcompute!55] (0,-2) rectangle ++(0.8,0.6); \node[font=\scriptsize, anchor=west] at (1,-1.7) {bf16 weights};
\draw[fill=fcomm!45] (7,-2) rectangle ++(0.8,0.6); \node[font=\scriptsize, anchor=west] at (8,-1.7) {bf16 gradients};
\draw[fill=fmemory!60] (14.5,-2) rectangle ++(0.8,0.6); \node[font=\scriptsize, anchor=west] at (15.5,-1.7) {fp32 master + Adam $m,v$};
\end{tikzpicture}
\caption{ZeRO stages on four GPUs. Bar heights are proportional to bytes per parameter (2, 2
and 12). Each stage shards one more kind of state so that every GPU keeps only its $1/N_d$
slice; ZeRO-3 shards everything.}
\label{fig:dist:zero}
\end{figure}

\paragraph{ZeRO-1 and ZeRO-2.} Each rank \emph{owns} a $1/N_d$ shard of the parameters and keeps
fp32 master weights and Adam moments only for that shard. After backward, gradients are
\emph{reduce-scattered} (each rank receives the summed gradient of its shard), the owner runs
Adam on its shard, and the updated parameters are \emph{all-gathered}. Since reduce-scatter +
all-gather moves exactly the bytes of one all-reduce, stages 1 and 2 cost the same
communication as plain DP while cutting memory by up to $4\times$ and $8\times$. Stage 2 also
frees the non-owned gradient slices immediately after their reduction.

Our teaching implementation (\cref{lst:dist:zero}) assigns whole parameter tensors to owners
and uses reduce/broadcast, because the gloo backend lacks reduce-scatter; production versions
flatten parameters into one buffer and shard it element-wise for perfect balance.

\begin{lstlisting}[caption={ZeRO-1/2 with parameter-level ownership (\code{llmsys/dist/zero.py}, abridged).},label={lst:dist:zero}]
class ZeroAdamW:
    def __init__(self, params, lr, stage=1, **kw):
        self.params = [p for p in params if p.requires_grad]
        self.owner = partition_parameters(self.params, dist.get_world_size())
        mine = [p for p, o in zip(self.params, self.owner) if o == dist.get_rank()]
        self.local = torch.optim.AdamW(mine, lr=lr, **kw)   # states only for my shard

    @torch.no_grad()
    def step(self):
        for p, o in zip(self.params, self.owner):
            if self.stage == 1:
                dist.all_reduce(p.grad); p.grad /= self.world
            else:                                  # stage 2: only the owner keeps it
                dist.reduce(p.grad, dst=o)
                if dist.get_rank() == o: p.grad /= self.world
                else: p.grad = None
        self.local.step()                          # Adam on owned parameters only
        for p, o in zip(self.params, self.owner):
            dist.broadcast(p.data, src=o)          # everyone gets updated weights
\end{lstlisting}

\paragraph{ZeRO-3 and FSDP.} Stage 3 also shards the bf16 parameters. The model is wrapped in
\emph{units} (typically one transformer block each). Just before a unit's forward, its
parameters are all-gathered from all ranks; after the forward they are freed. The backward
pass all-gathers them again, computes gradients, and reduce-scatters the gradients so each
rank keeps only its shard. Communication is $1.5\times$ that of DP ($3\Psi$ vs.\ $2\Psi$
bytes), but only one or two units are ever materialized, so per-GPU memory falls to
$16\Psi/N_d$ plus a unit or two of full parameters. PyTorch's FSDP (Zhao et al., 2023)
implements this design; FSDP2 shards each parameter individually as a \code{DTensor}, which
composes cleanly with tensor parallelism.

\begin{systemslens}[Making FSDP fast]
\begin{itemize}
\item \textbf{Prefetching.} Issue the all-gather for unit $i+1$ while unit $i$ computes
  (forward), and for unit $i-1$ during unit $i$'s backward. With enough compute per unit, the
  all-gathers are completely hidden.
\item \textbf{Hybrid sharding (HSDP).} Shard within a node (fast NVLink) and replicate across
  nodes, combining FSDP's memory savings with DP's lighter cross-node traffic (a plain
  all-reduce of gradient shards across replicas).
\item \textbf{Mixed-precision communication.} All-gather bf16 parameters and reduce-scatter
  bf16 (or fp32 for stability) gradients; keep fp32 master shards locally.
\item \textbf{Offload.} ZeRO-Offload/Infinity move optimizer states and even parameters to CPU
  memory or NVMe, trading PCIe bandwidth for capacity---useful for fine-tuning big models on
  few GPUs, rarely for pretraining throughput.
\end{itemize}
\end{systemslens}

\begin{figure}[t]
\centering
\begin{tikzpicture}[x=0.62cm, y=0.55cm]
\node[font=\scriptsize, anchor=east] at (-0.2,1.3) {compute};
\node[font=\scriptsize, anchor=east] at (-0.2,0.3) {comm};
\foreach \i in {1,2,3,4} {
  \draw[fill=fcompute!35] (\i*1.6-0.4,1) rectangle ++(1.4,0.6) node[midway, font=\tiny] {F\i};
  \draw[fill=fcomm!35] (\i*1.6-1.6,0) rectangle ++(1.1,0.6) node[midway, font=\tiny] {AG\i};
}
\foreach \i/\j in {4/0,3/1,2/2,1/3} {
  \draw[fill=fcompute!65] (7.2+\j*2.3,1) rectangle ++(2.1,0.6) node[midway, font=\tiny] {B\i};
  \draw[fill=fcomm!35] (6.0+\j*2.3,0) rectangle ++(1.1,0.6) node[midway, font=\tiny] {AG\i};
  \draw[fill=fmemory!50] (9.3+\j*2.3,-1) rectangle ++(1.1,0.6) node[midway, font=\tiny] {RS\i};
}
\node[font=\scriptsize, anchor=east] at (-0.2,-0.7) {comm 2};
\draw[-{Stealth}, srcgray] (0,-1.8) -- (19.5,-1.8) node[right, font=\scriptsize] {time};
\end{tikzpicture}
\caption{FSDP with prefetching for a 4-unit model. The all-gather (AG) of unit $i+1$ runs while
unit $i$ computes forward (F); in backward (B) units are re-gathered in reverse order and each
unit's gradients are reduce-scattered (RS) as soon as they exist. Only one or two units'
full parameters are ever resident.}
\label{fig:dist:fsdp}
\end{figure}

% ======================================================================
\section{Tensor parallelism}\label{sec:dist:tp}

ZeRO shards \emph{storage}, but every layer still runs on one device with full-size
activations. Tensor parallelism (Megatron-LM; Shoeybi et al., 2019) splits the
\emph{computation} of each layer. Take a linear layer $\mat Y=\mat X\mat W^\top$ with
$\mat W\in\R^{d_{\text{out}}\times d_{\text{in}}}$ and $t$ ranks.

\begin{description}
\item[Column-parallel] Split $\mat W$ by output rows, $\mat W=[\mat W_1;\dots;\mat W_t]$. Rank
  $r$ computes $\mat Y_r=\mat X\mat W_r^\top$, a slice of the output features. Input $\mat X$ must
  be replicated. In backward, each rank's $\partial\mathcal L/\partial\mat X$ is only the
  contribution of its slice, so the input gradient needs an \textbf{all-reduce}.
\item[Row-parallel] Split $\mat W$ by input columns. Rank $r$ consumes the matching slice
  $\mat X_r$ of the input and produces a partial result $\mat X_r\mat W_r^\top$; the true output is
  the sum over ranks, an \textbf{all-reduce} in the forward pass. Its backward needs no
  communication: the output gradient is already replicated.
\end{description}

The trick is to chain them. In the MLP $\mat Y=\GELU(\mat X\mat A^\top)\mat B^\top$, make $\mat A$
column-parallel and $\mat B$ row-parallel. Each rank then holds a slice of the hidden features,
applies GELU locally (element-wise, so no communication), and feeds that slice straight into
its row-parallel slice of $\mat B$. One all-reduce at the end of the forward pass and one at the
start of the backward pass suffice (\cref{fig:dist:tp}). Megatron calls the two conjugate
operators $f$ (identity forward, all-reduce backward) and $g$ (all-reduce forward, identity
backward). Self-attention splits the same way: heads are distributed column-parallel through
the Q/K/V projections (each rank runs complete heads, so softmax stays local) and the output
projection is row-parallel. A transformer layer thus needs \textbf{two all-reduces in forward
and two in backward}.

\begin{figure}[t]
\centering
\begin{tikzpicture}[
  op/.style={draw, rounded corners=1.5pt, minimum width=1.6cm, minimum height=0.7cm,
     font=\scriptsize, align=center, fill=#1},
  arr/.style={-{Stealth[length=1.6mm]}, thick}]
\node[op=fdata!20] (x) at (0,0) {$\mat X$\\ (replicated)};
\node[op=fcomm!35, minimum width=0.8cm] (f) at (1.9,0) {$f$};
\node[op=fcompute!25] (a0) at (3.8,0.75) {$\mat X\mat A_0^\top$};
\node[op=fcompute!25] (a1) at (3.8,-0.75) {$\mat X\mat A_1^\top$};
\node[op=fcompute!12] (g0) at (5.8,0.75) {GELU};
\node[op=fcompute!12] (g1) at (5.8,-0.75) {GELU};
\node[op=fcompute!25] (b0) at (7.8,0.75) {$\mat H_0\mat B_0^\top$};
\node[op=fcompute!25] (b1) at (7.8,-0.75) {$\mat H_1\mat B_1^\top$};
\node[op=fcomm!35, minimum width=0.8cm] (g) at (9.7,0) {$g$};
\node[op=fdata!20] (y) at (11.5,0) {$\mat Y$\\ (replicated)};
\draw[arr] (x) -- (f); \draw[arr] (f) -- (a0); \draw[arr] (f) -- (a1);
\draw[arr] (a0) -- (g0); \draw[arr] (a1) -- (g1); \draw[arr] (g0) -- (b0); \draw[arr] (g1) -- (b1);
\draw[arr] (b0) -- (g); \draw[arr] (b1) -- (g); \draw[arr] (g) -- (y);
\node[font=\scriptsize, text=srcgray] at (3.8,1.55) {column-parallel $\mat A$};
\node[font=\scriptsize, text=srcgray] at (7.8,1.55) {row-parallel $\mat B$};
\node[font=\scriptsize, align=center, fcomm] at (1.9,-1.35) {fwd: identity\\ bwd: all-reduce};
\node[font=\scriptsize, align=center, fcomm] at (9.7,-1.35) {fwd: all-reduce\\ bwd: identity};
\node[font=\scriptsize, text=srcgray] at (5.8,-2.2) {top row: rank 0 \quad bottom row: rank 1};
\end{tikzpicture}
\caption{Megatron-style tensor-parallel MLP on two ranks. The hidden activation $\mat H$ never
leaves its rank; only the $f$ and $g$ operators communicate.}
\label{fig:dist:tp}
\end{figure}

\begin{lstlisting}[caption={The $f$/$g$ operators and parallel linear layers (\code{llmsys/dist/tensor\_parallel.py}, abridged).},label={lst:dist:tp}]
class CopyToTP(torch.autograd.Function):          # f
    @staticmethod
    def forward(ctx, x): return x
    @staticmethod
    def backward(ctx, g):
        g = g.clone(); dist.all_reduce(g); return g

class ReduceFromTP(torch.autograd.Function):      # g
    @staticmethod
    def forward(ctx, x):
        x = x.clone(); dist.all_reduce(x); return x
    @staticmethod
    def backward(ctx, g): return g

class ColumnParallelLinear(nn.Module):            # holds out_features / t rows
    def forward(self, x):
        return F.linear(CopyToTP.apply(x), self.weight, self.bias)

class RowParallelLinear(nn.Module):               # holds in_features / t columns
    def forward(self, x_local):
        y = ReduceFromTP.apply(F.linear(x_local, self.weight))
        return y + self.bias                      # bias added once, after the sum
\end{lstlisting}

\begin{napkinmath}[Why TP stays inside a node]
A 70B layer ($h=8192$) with $s=4096$, $b=1$ and $t=8$: each all-reduce carries the
$s\cdot b\cdot h$ bf16 activation, \SI{67}{MB}; four per layer per step at
$2\cdot\frac78\cdot\SI{67}{MB}$ each $\approx\SI{0.47}{GB}$ per rank. At NVLink's
\SI{450}{GB/s}: \SI{1.0}{ms} per layer, versus $\approx\SI{5.4}{ms}$ of compute per layer per
GPU at 50\% MFU. Over \SI{50}{GB/s} IB the same traffic takes \SI{9.4}{ms}---longer than the
compute. TP is therefore almost always confined to the 8 GPUs of an NVLink domain (or the 72
of an NVL72 rack).
\end{napkinmath}

\paragraph{Hiding TP communication.} The all-reduces in \cref{fig:dist:tp} sit on the
critical path: the next layer cannot start until $g$ completes. Two families of techniques
overlap them anyway. \emph{Decomposition} (Wang et al., 2023; ``async TP'' in PyTorch, Flux,
and NVIDIA's userbuffers) splits the GEMM adjacent to a collective into chunks along the
sequence dimension, so that chunk $i$'s reduce-scatter or all-gather proceeds while chunk $i+1$
is multiplied. \emph{Micro-batch interleaving} runs two independent micro-batches through the
layer out of phase, so one computes while the other communicates (the approach behind
DeepSeek-V3's dual micro-batch overlap). Both turn the $\SI{1.0}{ms}$ of the napkin math
above into almost-free communication when the chunked GEMMs remain large enough to run
efficiently.

\begin{keyidea}[TP degree is a trade-off, not a free lunch]
Each doubling of $t$ halves every GEMM's size along one dimension. A $4096\times8192$ GEMM split
16 ways becomes $4096\times512$, which may no longer reach peak tensor-core efficiency, while the
activation all-reduce stays the same size. Choose the smallest $t$ that makes memory fit.
\end{keyidea}

\paragraph{Sequence parallelism.} TP leaves LayerNorm, dropout and the residual stream
replicated: every rank stores the full $s\times b\times h$ activations there. Korthikanti et al.\
(2022) split these regions along the \emph{sequence} dimension instead and replace each
all-reduce with a reduce-scatter (entering the sequence-parallel region) plus an all-gather
(leaving it). Communication volume is unchanged, since all-reduce = reduce-scatter +
all-gather, while the replicated activations shrink by $t$.

\paragraph{Vocabulary parallelism.} The embedding table and output head ($V\times h$, up to
1.2B parameters for $V=152$k, $h=8192$) are split along the vocabulary. The cross-entropy is then
computed without gathering logits: each rank computes its local max and local
$\sum\exp$ over its vocabulary slice, two small all-reduces produce the global normalizer, and
the target logit comes from whichever rank owns the target ID.

% ======================================================================
\section{Pipeline parallelism}\label{sec:dist:pp}

Pipeline parallelism assigns contiguous groups of layers (stages) to different devices. Only
the activations at stage boundaries cross devices---$s\cdot b\cdot h$ values per micro-batch,
point-to-point---so PP tolerates slow inter-node links well. The catch is idleness: stage 2
cannot start until stage 1 finishes. GPipe (Huang et al., 2019) splits each batch into $m$
micro-batches and streams them through the stages: all forwards, then all backwards, then one
synchronous optimizer step. With $p$ stages, per-micro-batch forward time $t_f$ and backward
time $t_b$, each stage is idle for $(p-1)(t_f+t_b)$ out of $(m+p-1)(t_f+t_b)$:
\begin{equation}
\text{bubble fraction} = \frac{p-1}{m+p-1}.
\label{eq:dist:bubble}
\end{equation}
To keep the bubble small one needs $m\gg p$ (e.g.\ $m\ge4p$ gives $<20\%$). But GPipe holds the
activations of \emph{all} $m$ micro-batches at once before the first backward.

\paragraph{1F1B.} PipeDream-Flush / Megatron's \emph{one-forward-one-backward} schedule (Narayanan
et al., 2021) has the same bubble but starts each micro-batch's backward as early as possible:
stage $s$ runs $p-s-1$ warm-up forwards, then alternates one forward and one backward, then
drains. A stage never holds more than $p-s$ micro-batches of activations, independent of $m$.
\Cref{fig:dist:pipe} shows both schedules as produced by \code{llmsys.dist.pipeline.simulate}.

\begin{figure}[t]
\centering
\begin{tikzpicture}
\node[font=\footnotesize\bfseries, anchor=west] at (-1.2,0.45) {GPipe ($p=4$, $m=6$): peak in-flight micro-batches per stage = 6, 6, 6, 6};
\draw[fill=fcompute!35] (0.000,0.000) rectangle (0.360,-0.425) node[midway,font=\tiny] {1};
\draw[fill=fcompute!35] (0.360,-0.500) rectangle (0.720,-0.925) node[midway,font=\tiny] {1};
\draw[fill=fcompute!35] (0.720,-1.000) rectangle (1.080,-1.425) node[midway,font=\tiny] {1};
\draw[fill=fcompute!35] (1.080,-1.500) rectangle (1.440,-1.925) node[midway,font=\tiny] {1};
\draw[fill=fcompute!35] (0.360,0.000) rectangle (0.720,-0.425) node[midway,font=\tiny] {2};
\draw[fill=fcompute!35] (0.720,-0.500) rectangle (1.080,-0.925) node[midway,font=\tiny] {2};
\draw[fill=fcompute!35] (1.080,-1.000) rectangle (1.440,-1.425) node[midway,font=\tiny] {2};
\draw[fill=fcompute!35] (1.440,-1.500) rectangle (1.800,-1.925) node[midway,font=\tiny] {2};
\draw[fill=fcompute!35] (0.720,0.000) rectangle (1.080,-0.425) node[midway,font=\tiny] {3};
\draw[fill=fcompute!35] (1.080,-0.500) rectangle (1.440,-0.925) node[midway,font=\tiny] {3};
\draw[fill=fcompute!35] (1.440,-1.000) rectangle (1.800,-1.425) node[midway,font=\tiny] {3};
\draw[fill=fcompute!35] (1.800,-1.500) rectangle (2.160,-1.925) node[midway,font=\tiny] {3};
\draw[fill=fcompute!35] (1.080,0.000) rectangle (1.440,-0.425) node[midway,font=\tiny] {4};
\draw[fill=fcompute!35] (1.440,-0.500) rectangle (1.800,-0.925) node[midway,font=\tiny] {4};
\draw[fill=fcompute!35] (1.800,-1.000) rectangle (2.160,-1.425) node[midway,font=\tiny] {4};
\draw[fill=fcompute!35] (2.160,-1.500) rectangle (2.520,-1.925) node[midway,font=\tiny] {4};
\draw[fill=fcompute!35] (1.440,0.000) rectangle (1.800,-0.425) node[midway,font=\tiny] {5};
\draw[fill=fcompute!35] (1.800,-0.500) rectangle (2.160,-0.925) node[midway,font=\tiny] {5};
\draw[fill=fcompute!35] (2.160,-1.000) rectangle (2.520,-1.425) node[midway,font=\tiny] {5};
\draw[fill=fcompute!35] (2.520,-1.500) rectangle (2.880,-1.925) node[midway,font=\tiny] {5};
\draw[fill=fcompute!35] (1.800,0.000) rectangle (2.160,-0.425) node[midway,font=\tiny] {6};
\draw[fill=fcompute!35] (2.160,-0.500) rectangle (2.520,-0.925) node[midway,font=\tiny] {6};
\draw[fill=fcompute!35] (2.520,-1.000) rectangle (2.880,-1.425) node[midway,font=\tiny] {6};
\draw[fill=fcompute!35] (2.880,-1.500) rectangle (3.240,-1.925) node[midway,font=\tiny] {6};
\draw[fill=fcomm!35] (3.240,-1.500) rectangle (3.960,-1.925) node[midway,font=\tiny] {6};
\draw[fill=fcomm!35] (3.960,-1.000) rectangle (4.680,-1.425) node[midway,font=\tiny] {6};
\draw[fill=fcomm!35] (3.960,-1.500) rectangle (4.680,-1.925) node[midway,font=\tiny] {5};
\draw[fill=fcomm!35] (4.680,-0.500) rectangle (5.400,-0.925) node[midway,font=\tiny] {6};
\draw[fill=fcomm!35] (4.680,-1.000) rectangle (5.400,-1.425) node[midway,font=\tiny] {5};
\draw[fill=fcomm!35] (4.680,-1.500) rectangle (5.400,-1.925) node[midway,font=\tiny] {4};
\draw[fill=fcomm!35] (5.400,0.000) rectangle (6.120,-0.425) node[midway,font=\tiny] {6};
\draw[fill=fcomm!35] (5.400,-0.500) rectangle (6.120,-0.925) node[midway,font=\tiny] {5};
\draw[fill=fcomm!35] (5.400,-1.000) rectangle (6.120,-1.425) node[midway,font=\tiny] {4};
\draw[fill=fcomm!35] (5.400,-1.500) rectangle (6.120,-1.925) node[midway,font=\tiny] {3};
\draw[fill=fcomm!35] (6.120,0.000) rectangle (6.840,-0.425) node[midway,font=\tiny] {5};
\draw[fill=fcomm!35] (6.120,-0.500) rectangle (6.840,-0.925) node[midway,font=\tiny] {4};
\draw[fill=fcomm!35] (6.120,-1.000) rectangle (6.840,-1.425) node[midway,font=\tiny] {3};
\draw[fill=fcomm!35] (6.120,-1.500) rectangle (6.840,-1.925) node[midway,font=\tiny] {2};
\draw[fill=fcomm!35] (6.840,0.000) rectangle (7.560,-0.425) node[midway,font=\tiny] {4};
\draw[fill=fcomm!35] (6.840,-0.500) rectangle (7.560,-0.925) node[midway,font=\tiny] {3};
\draw[fill=fcomm!35] (6.840,-1.000) rectangle (7.560,-1.425) node[midway,font=\tiny] {2};
\draw[fill=fcomm!35] (6.840,-1.500) rectangle (7.560,-1.925) node[midway,font=\tiny] {1};
\draw[fill=fcomm!35] (7.560,0.000) rectangle (8.280,-0.425) node[midway,font=\tiny] {3};
\draw[fill=fcomm!35] (7.560,-0.500) rectangle (8.280,-0.925) node[midway,font=\tiny] {2};
\draw[fill=fcomm!35] (7.560,-1.000) rectangle (8.280,-1.425) node[midway,font=\tiny] {1};
\draw[fill=fcomm!35] (8.280,0.000) rectangle (9.000,-0.425) node[midway,font=\tiny] {2};
\draw[fill=fcomm!35] (8.280,-0.500) rectangle (9.000,-0.925) node[midway,font=\tiny] {1};
\draw[fill=fcomm!35] (9.000,0.000) rectangle (9.720,-0.425) node[midway,font=\tiny] {1};
\node[anchor=east,font=\scriptsize] at (-0.1,-0.210) {GPU 0};
\node[anchor=east,font=\scriptsize] at (-0.1,-0.710) {GPU 1};
\node[anchor=east,font=\scriptsize] at (-0.1,-1.210) {GPU 2};
\node[anchor=east,font=\scriptsize] at (-0.1,-1.710) {GPU 3};
\begin{scope}[yshift=-2.9cm]
\node[font=\footnotesize\bfseries, anchor=west] at (-1.2,0.45) {1F1B ($p=4$, $m=6$): peak in-flight micro-batches per stage = 4, 3, 2, 1};
\draw[fill=fcompute!35] (0.000,0.000) rectangle (0.360,-0.425) node[midway,font=\tiny] {1};
\draw[fill=fcompute!35] (0.360,-0.500) rectangle (0.720,-0.925) node[midway,font=\tiny] {1};
\draw[fill=fcompute!35] (0.720,-1.000) rectangle (1.080,-1.425) node[midway,font=\tiny] {1};
\draw[fill=fcompute!35] (1.080,-1.500) rectangle (1.440,-1.925) node[midway,font=\tiny] {1};
\draw[fill=fcompute!35] (0.360,0.000) rectangle (0.720,-0.425) node[midway,font=\tiny] {2};
\draw[fill=fcompute!35] (0.720,-0.500) rectangle (1.080,-0.925) node[midway,font=\tiny] {2};
\draw[fill=fcompute!35] (1.080,-1.000) rectangle (1.440,-1.425) node[midway,font=\tiny] {2};
\draw[fill=fcomm!35] (1.440,-1.500) rectangle (2.160,-1.925) node[midway,font=\tiny] {1};
\draw[fill=fcompute!35] (0.720,0.000) rectangle (1.080,-0.425) node[midway,font=\tiny] {3};
\draw[fill=fcompute!35] (1.080,-0.500) rectangle (1.440,-0.925) node[midway,font=\tiny] {3};
\draw[fill=fcomm!35] (2.160,-1.000) rectangle (2.880,-1.425) node[midway,font=\tiny] {1};
\draw[fill=fcompute!35] (2.160,-1.500) rectangle (2.520,-1.925) node[midway,font=\tiny] {2};
\draw[fill=fcompute!35] (1.080,0.000) rectangle (1.440,-0.425) node[midway,font=\tiny] {4};
\draw[fill=fcomm!35] (2.880,-0.500) rectangle (3.600,-0.925) node[midway,font=\tiny] {1};
\draw[fill=fcompute!35] (2.880,-1.000) rectangle (3.240,-1.425) node[midway,font=\tiny] {3};
\draw[fill=fcomm!35] (2.520,-1.500) rectangle (3.240,-1.925) node[midway,font=\tiny] {2};
\draw[fill=fcomm!35] (3.600,0.000) rectangle (4.320,-0.425) node[midway,font=\tiny] {1};
\draw[fill=fcompute!35] (3.600,-0.500) rectangle (3.960,-0.925) node[midway,font=\tiny] {4};
\draw[fill=fcomm!35] (3.240,-1.000) rectangle (3.960,-1.425) node[midway,font=\tiny] {2};
\draw[fill=fcompute!35] (3.240,-1.500) rectangle (3.600,-1.925) node[midway,font=\tiny] {3};
\draw[fill=fcompute!35] (4.320,0.000) rectangle (4.680,-0.425) node[midway,font=\tiny] {5};
\draw[fill=fcomm!35] (3.960,-0.500) rectangle (4.680,-0.925) node[midway,font=\tiny] {2};
\draw[fill=fcompute!35] (3.960,-1.000) rectangle (4.320,-1.425) node[midway,font=\tiny] {4};
\draw[fill=fcomm!35] (3.600,-1.500) rectangle (4.320,-1.925) node[midway,font=\tiny] {3};
\draw[fill=fcomm!35] (4.680,0.000) rectangle (5.400,-0.425) node[midway,font=\tiny] {2};
\draw[fill=fcompute!35] (4.680,-0.500) rectangle (5.040,-0.925) node[midway,font=\tiny] {5};
\draw[fill=fcomm!35] (4.320,-1.000) rectangle (5.040,-1.425) node[midway,font=\tiny] {3};
\draw[fill=fcompute!35] (4.320,-1.500) rectangle (4.680,-1.925) node[midway,font=\tiny] {4};
\draw[fill=fcompute!35] (5.400,0.000) rectangle (5.760,-0.425) node[midway,font=\tiny] {6};
\draw[fill=fcomm!35] (5.040,-0.500) rectangle (5.760,-0.925) node[midway,font=\tiny] {3};
\draw[fill=fcompute!35] (5.040,-1.000) rectangle (5.400,-1.425) node[midway,font=\tiny] {5};
\draw[fill=fcomm!35] (4.680,-1.500) rectangle (5.400,-1.925) node[midway,font=\tiny] {4};
\draw[fill=fcomm!35] (5.760,0.000) rectangle (6.480,-0.425) node[midway,font=\tiny] {3};
\draw[fill=fcompute!35] (5.760,-0.500) rectangle (6.120,-0.925) node[midway,font=\tiny] {6};
\draw[fill=fcomm!35] (5.400,-1.000) rectangle (6.120,-1.425) node[midway,font=\tiny] {4};
\draw[fill=fcompute!35] (5.400,-1.500) rectangle (5.760,-1.925) node[midway,font=\tiny] {5};
\draw[fill=fcomm!35] (6.120,-0.500) rectangle (6.840,-0.925) node[midway,font=\tiny] {4};
\draw[fill=fcompute!35] (6.120,-1.000) rectangle (6.480,-1.425) node[midway,font=\tiny] {6};
\draw[fill=fcomm!35] (5.760,-1.500) rectangle (6.480,-1.925) node[midway,font=\tiny] {5};
\draw[fill=fcomm!35] (6.840,0.000) rectangle (7.560,-0.425) node[midway,font=\tiny] {4};
\draw[fill=fcomm!35] (6.480,-1.000) rectangle (7.200,-1.425) node[midway,font=\tiny] {5};
\draw[fill=fcompute!35] (6.480,-1.500) rectangle (6.840,-1.925) node[midway,font=\tiny] {6};
\draw[fill=fcomm!35] (7.200,-0.500) rectangle (7.920,-0.925) node[midway,font=\tiny] {5};
\draw[fill=fcomm!35] (6.840,-1.500) rectangle (7.560,-1.925) node[midway,font=\tiny] {6};
\draw[fill=fcomm!35] (7.920,0.000) rectangle (8.640,-0.425) node[midway,font=\tiny] {5};
\draw[fill=fcomm!35] (7.560,-1.000) rectangle (8.280,-1.425) node[midway,font=\tiny] {6};
\draw[fill=fcomm!35] (8.280,-0.500) rectangle (9.000,-0.925) node[midway,font=\tiny] {6};
\draw[fill=fcomm!35] (9.000,0.000) rectangle (9.720,-0.425) node[midway,font=\tiny] {6};
\node[anchor=east,font=\scriptsize] at (-0.1,-0.210) {GPU 0};
\node[anchor=east,font=\scriptsize] at (-0.1,-0.710) {GPU 1};
\node[anchor=east,font=\scriptsize] at (-0.1,-1.210) {GPU 2};
\node[anchor=east,font=\scriptsize] at (-0.1,-1.710) {GPU 3};
\end{scope}
\draw[fill=fcompute!35] (0,-5.4) rectangle ++(0.4,0.3); \node[font=\scriptsize, anchor=west] at (0.45,-5.25) {forward ($t_f=1$)};
\draw[fill=fcomm!35] (3.2,-5.4) rectangle ++(0.8,0.3); \node[font=\scriptsize, anchor=west] at (4.05,-5.25) {backward ($t_b=2$); number = micro-batch};
\end{tikzpicture}
\caption{Pipeline schedules simulated with \code{llmsys.dist.pipeline}. Both take 27 time
units and idle for $(p-1)/(m+p-1)=1/3$ of the time (white gaps), but 1F1B frees each
micro-batch's activations as soon as its backward completes.}
\label{fig:dist:pipe}
\end{figure}

\paragraph{Shrinking the bubble.} \emph{Interleaved} 1F1B gives each device $v$ non-contiguous
``virtual'' stages (e.g.\ layers 1--4 and 17--20), shrinking the bubble by $v$ at the cost of
$v$ times more point-to-point messages. \emph{Zero-bubble} schedules (Qi et al., 2024) split the
backward pass into the input-gradient part $B$ (on the critical path) and the weight-gradient
part $W$ (which can be deferred), and fill the bubbles with $W$ work. DeepSeek-V3's
\emph{DualPipe} feeds micro-batches from both ends of the pipeline and overlaps
computation with the all-to-all communication of its MoE layers.

\begin{systemslens}[Practical pipeline concerns]
\begin{itemize}
\item \textbf{Balance.} The slowest stage sets the pace. The embedding layer and the output
  head with its large softmax make the first and last stages heavier; Megatron lets you assign
  fewer transformer layers to them.
\item \textbf{Micro-batch count vs.\ global batch.} $m$ is bounded by global batch $/$ (DP
  degree $\times$ micro-batch size). With a fixed global batch, more DP means fewer micro-batches
  per pipeline and a bigger bubble---a central tension of 3D parallelism.
\item \textbf{Activation memory is uneven}: stage 0 holds $p$ micro-batches under 1F1B, the last
  stage only one.
\item \textbf{Weight updates are synchronous} (flush at the end of each batch), which keeps the
  math identical to non-pipelined training; asynchronous PipeDream variants with weight
  stashing are rarely used for LLMs.
\end{itemize}
\end{systemslens}

% ======================================================================
\section{Sequence and context parallelism for long contexts}\label{sec:dist:cp}

At 128k-token contexts the activations of even one sequence overwhelm a GPU, and attention's
$O(s^2)$ compute dominates. \emph{Context parallelism} (CP) splits the sequence of every
layer---not just the LayerNorm regions---across $c$ devices. The MLP and projections are
per-token and need no communication; attention needs every query to see every (earlier) key.

\paragraph{Ring attention} (Liu et al., 2023). Each device keeps its query chunk and passes its
key/value chunk around a ring. At each of $c$ steps it computes attention of its queries
against the visiting K/V chunk and folds the result into its output with the online-softmax
rescaling of FlashAttention (\cref{ch:attention}): keep a running max $m$ and normalizer
$\ell$ per query, and rescale the partial output when a new block raises the max. Sending the
next K/V chunk overlaps with computing on the current one. With causal masking, naive
contiguous chunking makes the last device do $c$ times more work than the first;
\emph{zig-zag} assignment (device $i$ gets chunks $i$ and $2c-1-i$) balances it.

\begin{figure}[t]
\centering
\begin{tikzpicture}[x=0.55cm, y=0.55cm]
\foreach \st in {0,1,2,3} {
  \begin{scope}[xshift=\st*3.9cm]
  \node[font=\scriptsize] at (2,4.7) {step \st};
  \foreach \q in {0,1,2,3} {
    \pgfmathtruncatemacro{\kv}{mod(\q-\st+4,4)}
    \foreach \k in {0,1,2,3} {
      \pgfmathtruncatemacro{\active}{\k==\kv ? 1 : 0}
      \pgfmathtruncatemacro{\masked}{\k>\q ? 1 : 0}
      \ifnum\masked=1 \def\col{white}\else
        \ifnum\active=1 \def\col{fcompute!60}\else\def\col{srcgray!12}\fi\fi
      \draw[fill=\col] (\k,3-\q) rectangle ++(1,1);
    }
    \node[font=\tiny, anchor=east] at (0,3.5-\q) {D\q};
  }
  \node[font=\tiny] at (2,-0.4) {K/V blocks};
  \end{scope}
}
\end{tikzpicture}
\caption{Ring attention on four devices with contiguous chunks. Row $q$ is device $q$'s query
block; at step $t$ it holds K/V block $(q-t)\bmod 4$ (dark). White cells are causally masked:
device 0 has useful work only at step 0 while device 3 works at every step---the imbalance that
zig-zag chunk assignment removes.}
\label{fig:dist:ringattn}
\end{figure}

\paragraph{Ulysses} (Jacobs et al., 2023). Alternatively, an all-to-all before attention
re-partitions from ``all heads, $1/c$ of the sequence'' to ``$1/c$ of the heads, the full
sequence'', runs ordinary (Flash)attention locally, and a second all-to-all switches back.
Communication per device is $O(s\,h/c)$ and independent of the number of steps, but the degree
is limited by the number of (KV) heads.

% ======================================================================
\section{Expert parallelism}\label{sec:dist:ep}

Mixture-of-experts layers (\cref{ch:modern}) replace the MLP with $E$ expert MLPs and a router
that sends each token to its top-$k$ experts. With expert parallelism, experts are spread over
$e$ devices. Each MoE layer then needs two all-to-alls: \emph{dispatch} (send each token's hidden
state to the devices hosting its chosen experts) and \emph{combine} (send results back).
Per token and layer, about $2k\cdot h$ values cross the network in each direction. Systems
issues include load imbalance (a popular expert becomes a straggler; auxiliary balancing losses
or bias-based balancing mitigate it), token dropping when an expert exceeds its \emph{capacity}
($\text{capacity factor}\times$ tokens$/E$), and topology-aware routing that limits each token to
experts on a few nodes so that most all-to-all traffic stays on NVLink (as in DeepSeek-V3).

\begin{napkinmath}[All-to-all volume of a large MoE]
Take a DeepSeek-V3-like layer ($h=7168$, top-$k=8$ routed experts) and 8{,}192 tokens per GPU.
Dispatch sends $8192\cdot8\cdot7168$ bf16 values $\approx\SI{0.94}{GB}$ per GPU per MoE layer, and
combine sends the same back---every layer, forward and backward. At \SI{50}{GB/s} that is
$\approx\SI{19}{ms}$ per direction per layer, comparable to the expert compute itself. Hence the
emphasis on node-limited routing (most traffic on NVLink), FP8 dispatch, and schedules that
overlap one micro-batch's all-to-all with another's computation.
\end{napkinmath}

% ======================================================================
\section{Putting it together: 3D and 4D parallelism}\label{sec:dist:3d}

Large runs compose these strategies on a \emph{device mesh}. The ordering follows
communication intensity: TP (and SP) innermost on NVLink; CP next; PP across nodes; DP/FSDP
outermost (\cref{fig:dist:mesh}). In PyTorch, \code{init\_device\_mesh("cuda", (dp, pp, tp))}
creates the sub-groups and FSDP2/TP/PP APIs each take one mesh dimension.

\begin{figure}[t]
\centering
\begin{tikzpicture}[x=0.78cm, y=0.6cm]
\foreach \node in {0,1,2,3} {
  \pgfmathsetmacro{\X}{\node*4.6}
  \draw[rounded corners=2pt, fill=srcgray!8, draw=srcgray] (\X-0.3,-0.3) rectangle ++(4.0,4.6);
  \node[font=\tiny] at (\X+1.7,4.6) {node \node};
  \foreach \g in {0,...,7} {
    \pgfmathsetmacro{\gx}{\X + mod(\g,4)*0.85}
    \pgfmathsetmacro{\gy}{3.2 - floor(\g/4)*1.1}
    \pgfmathtruncatemacro{\rank}{\node*8+\g}
    \pgfmathtruncatemacro{\pp}{mod(\node,2)}
    \ifnum\pp=0 \def\col{fcompute!30}\else\def\col{fcompute!60}\fi
    \draw[fill=\col] (\gx,\gy) rectangle ++(0.7,0.8);
    \node[font=\tiny] at (\gx+0.35,\gy+0.4) {\rank};
  }
  \draw[decorate, decoration={brace, amplitude=3pt}, fcomm] (\X,2.0) -- (\X+3.25,2.0)
    node[midway, below=2pt, font=\tiny, fcomm, fill=white, inner sep=1pt] {TP group};
}
\draw[-{Stealth}, thick, napkin] (3.7,0.6) -- (4.3,0.6); \node[font=\tiny, napkin] at (4.0,0.05) {PP};
\draw[-{Stealth}, thick, napkin] (12.9,0.6) -- (13.5,0.6); \node[font=\tiny, napkin] at (13.2,0.05) {PP};
\draw[{Stealth}-{Stealth}, thick, keygreen] (1.7,-0.8) -- (10.9,-0.8) node[midway, below, font=\tiny, keygreen] {DP replica 0 $\leftrightarrow$ DP replica 1 (gradient reduce over IB)};
\end{tikzpicture}
\caption{A 32-GPU mesh with TP $=8$ (one node), PP $=2$ (light/dark nodes are stages 0 and 1)
and DP $=2$. Rank $=$ dp$\cdot16$ + pp$\cdot8$ + tp: TP neighbors are adjacent ranks on the same
NVSwitch; PP sends activations to the next node; DP all-reduces span the whole cluster.}
\label{fig:dist:mesh}
\end{figure}

\begin{table}[t]
\centering\footnotesize
\begin{tabularx}{\linewidth}{lXXXl}
\toprule
strategy & splits & saves memory of & communication & typical link \\
\midrule
DP & batch & --- & all-reduce grads, $2\Psi$/step & IB \\
ZeRO-1/2 & batch & optimizer (+grads) & same as DP & IB \\
FSDP/ZeRO-3 & batch & everything & $3\Psi$/step, per-unit gathers & IB / NVLink \\
TP (+SP) & matrices & weights, activations & 4 all-reduces of $sbh$ per layer & NVLink \\
PP & layers & weights, activations & p2p $sbh$ per micro-batch & IB \\
CP & sequence & activations & K/V ring or all-to-all & NVLink / IB \\
EP & experts & expert weights & 2 all-to-alls per MoE layer & NVLink / IB \\
\bottomrule
\end{tabularx}
\caption{Parallelism strategies at a glance.}
\label{tab:dist:summary}
\end{table}

\paragraph{A planning recipe.}
\begin{enumerate}
\item Pick the smallest TP degree (within a node) that, together with PP and ZeRO, makes model
  states and activations fit; larger TP shrinks per-GPU matmuls and lowers efficiency.
\item Add PP when the model does not fit within one NVLink domain; choose $p$ so that layers
  divide evenly and $m\ge4p$ micro-batches are available.
\item Add CP only when sequence length requires it.
\item Spend all remaining GPUs on DP (with ZeRO-1 or FSDP), and check that the global batch
  stays below the critical batch size.
\item Verify with napkin math: memory per GPU, bubble fraction, exposed communication.
\end{enumerate}

\begin{napkinmath}[A 70B run on 1{,}024 H100s]
Choose TP $=8$, PP $=4$, DP $=32$ with ZeRO-1, $s=4096$, global batch 4M tokens.
\begin{itemize}
\item Parameters per GPU: $70\text{B}/(8\cdot4)=2.19$B. Model states with ZeRO-1:
  $(2+2+12/32)\times2.19\text{B}\approx\SI{9.6}{GB}$.
\item Activations: \cref{eq:pretrain:act} with FlashAttention gives \SI{1.14}{GB} per layer per
  sequence; TP+SP divides by 8; 20 layers per stage; stage 0 holds up to $p=4$ micro-batches
  under 1F1B: $1.14/8\times20\times4\approx\SI{11.4}{GB}$. Total $\approx\SI{21}{GB}$---comfortable on
  80\,GB, leaving room for larger micro-batches or less recomputation.
\item Bubble: 1{,}024 sequences per step / DP 32 = 32 micro-batches per pipeline;
  $(4-1)/(32+3)=8.6\%$. With PP $=8$ it would be 17.9\% (or 9.9\% interleaved with $v=2$).
\item DP gradient traffic: each rank reduces its \SI{4.4}{GB} bf16 shard over 32 replicas:
  $\approx\SI{0.17}{s}$ at \SI{50}{GB/s}, hidden behind $\approx\SI{4.1}{s}$ of compute per step at
  40\% MFU.
\end{itemize}
\end{napkinmath}

% ======================================================================
\section{From scratch to production APIs}\label{sec:dist:apis}

Everything above has a production counterpart in PyTorch. A training script is launched with
one process per GPU by \code{torchrun}, which sets \code{RANK}, \code{LOCAL\_RANK},
\code{WORLD\_SIZE} and the rendezvous address:

\begin{lstlisting}[style=shell]
torchrun --nnodes 4 --nproc-per-node 8 --rdzv-backend c10d \
         --rdzv-endpoint $HEAD_NODE:29500 train.py --tp 8 --dp 4
\end{lstlisting}

\begin{lstlisting}[caption={Composing TP and FSDP2 on a 2-D device mesh (sketch; PyTorch 2.4+ APIs).},label={lst:dist:prod}]
import os, torch, torch.distributed as dist
from torch.distributed.device_mesh import init_device_mesh
from torch.distributed.fsdp import fully_shard, MixedPrecisionPolicy
from torch.distributed.tensor.parallel import (parallelize_module,
    ColwiseParallel, RowwiseParallel)

torch.cuda.set_device(int(os.environ["LOCAL_RANK"]))
dist.init_process_group("nccl")
mesh = init_device_mesh("cuda", (4, 8), mesh_dim_names=("dp", "tp"))

model = build_model()                               # e.g. a Qwen3Model on "meta"
for blk in model.trf_blocks:                        # Megatron-style TP inside a node
    parallelize_module(blk, mesh["tp"], {
        "att.W_query": ColwiseParallel(), "att.W_key": ColwiseParallel(),
        "att.W_value": ColwiseParallel(), "att.out_proj": RowwiseParallel(),
        "ff.fc1": ColwiseParallel(), "ff.fc2": ColwiseParallel(),
        "ff.fc3": RowwiseParallel()})
mp = MixedPrecisionPolicy(param_dtype=torch.bfloat16, reduce_dtype=torch.float32)
for blk in model.trf_blocks:                        # FSDP units = blocks, over dp
    fully_shard(blk, mesh=mesh["dp"], mp_policy=mp)
fully_shard(model, mesh=mesh["dp"], mp_policy=mp)
opt = torch.optim.AdamW(model.parameters(), lr=3e-4, fused=True)
\end{lstlisting}

The mapping to this chapter is direct: \code{ColwiseParallel}/\code{RowwiseParallel} are our
column/row layers (the attention module must also be told that each rank now holds
$H/t$ heads); \code{fully\_shard} is ZeRO-3 with per-block units, prefetching and
reduce-scatter; the mesh encodes the rank layout of \cref{fig:dist:mesh}. Pipeline parallelism
is available through \code{torch.distributed.pipelining} (schedules \code{ScheduleGPipe},
\code{Schedule1F1B}, \code{ScheduleInterleaved1F1B}), and \code{torch.distributed.checkpoint}
saves and reshards DTensor-based state. Megatron-LM/Megatron-Core, DeepSpeed and TorchTitan are
the reference systems that package these pieces for LLM pretraining.

\begin{pitfall}[Test the small configuration first]
Every parallel layout must reproduce the single-device loss curve. Before a large launch, run
the same seed on 1 GPU and on a tiny mesh (e.g.\ $\text{TP}=2$, $\text{DP}=2$) for a few hundred
steps and compare losses to several decimal places. This is exactly what Lab~14's tests do, in
miniature.
\end{pitfall}

% ======================================================================
\section{Reliability at scale}\label{sec:dist:reliability}

At thousands of GPUs, failure is the steady state. Meta reported 466 job interruptions during
a 54-day period of Llama~3 405B pretraining on 16k H100s, 419 of them unexpected, most
attributed to GPU or HBM faults and other hardware issues (Llama Team, 2024)---roughly one every
three hours. Training infrastructure therefore treats recovery as a first-class path:
\begin{itemize}
\item \textbf{Fast, frequent checkpoints} (asynchronous, sharded, \cref{sec:pretrain:save}) and
  fast restarts (pre-warmed spare nodes, cached container images, quick NCCL
  re-initialization).
\item \textbf{Elastic launch.} \code{torchrun} restarts the worker group on failure; with
  resharding-capable checkpoints the job can resume on fewer or more nodes.
\item \textbf{Hang and straggler detection.} A single slow GPU (thermal throttling, a flaky
  link) slows every synchronous collective. Collective timeouts, per-rank step-time telemetry
  and NCCL flight recorders identify the culprit.
\item \textbf{Silent data corruption.} Rare hardware faults flip bits without crashing.
  Symptoms are loss spikes or NaNs on one replica; defenses include comparing gradient norms
  across DP replicas (they should match exactly after the all-reduce) and periodic
  known-answer tests.
\end{itemize}

\begin{pitfall}[Determinism across ranks]
DP correctness assumes every replica applies \emph{bitwise identical} updates. Non-deterministic
kernels or different reduction orders across ranks make replicas drift apart silently. Broadcast
parameters at start, use the same seeds for anything that must agree (dropout masks are fine to
differ; parameter initialization is not), and periodically check a weight checksum across ranks.
\end{pitfall}

% ======================================================================
\section{Summary}\label{sec:dist:summary}
\begin{itemize}
\item Collectives are the vocabulary of distributed training; the $\alpha$--$\beta$ model with the
  ring identity all-reduce = reduce-scatter + all-gather explains most of their costs.
\item DP scales throughput but not memory; bucketing and overlap hide gradient communication.
\item ZeRO-1/2 shard optimizer states and gradients at no extra communication; ZeRO-3/FSDP shard
  everything for 1.5$\times$ the traffic.
\item TP splits matrices with paired column/row layers (two all-reduces per pass per sub-block)
  and needs NVLink-class bandwidth; SP removes the remaining replicated activations.
\item PP splits layers; its bubble is $(p-1)/(m+p-1)$; 1F1B bounds activation memory, interleaving
  and zero-bubble schedules shrink idle time.
\item CP (ring attention, Ulysses) handles long sequences; EP distributes MoE experts with
  all-to-all.
\item Real jobs compose these on a device mesh ordered by bandwidth needs and must be engineered
  for frequent failure.
\end{itemize}

% ======================================================================
\section{Exercises}\label{sec:dist:exercises}

\begin{exercise}[Ring all-reduce cost]
Derive \cref{eq:dist:ring} by counting the bytes each rank sends in each of the $2(p-1)$ steps.
Then show that the recursive halving/doubling algorithm achieves the same bandwidth term with
only $2\log_2p$ latency terms when $p$ is a power of two.
\end{exercise}

\begin{exercise}[Two-level all-reduce]
A cluster has $N$ nodes of 8 GPUs, with intra-node bandwidth $B_1$ and per-GPU inter-node
bandwidth $B_2\ll B_1$. Design a hierarchical all-reduce (reduce-scatter in the node, all-reduce
across nodes on $1/8$ shards, all-gather in the node) and give its cost. Compare with a flat ring
over all $8N$ GPUs for $B_1=450$, $B_2=50$\,GB/s.
\end{exercise}

\begin{exercise}[Bucket size]
With $\alpha=\SI{20}{\micro s}$ and \SI{50}{GB/s}, what bucket size keeps the latency overhead
below 5\% of the transfer time? Why might a very large bucket still be bad for overlap?
\end{exercise}

\begin{exercise}[No-sync accumulation]
Modify \code{BucketedDDP} to support a \code{no\_sync()} context manager. Measure, with the
gloo backend, the time for 4-step gradient accumulation with and without it.
\end{exercise}

\begin{exercise}[ZeRO arithmetic]
Reproduce \cref{tab:dist:zero}. For a 13B model on 16 GPUs with 80\,GB each, which is the lowest
ZeRO stage that leaves at least 30\,GB per GPU for activations?
\end{exercise}

\begin{exercise}[TP backward]
Show that for a column-parallel linear layer the gradient with respect to its input is
$\sum_r \frac{\partial\mathcal L}{\partial \mat Y_r}\mat W_r$, which justifies the all-reduce in the
backward of $f$. Why does the row-parallel layer need no communication in backward?
\end{exercise}

\begin{exercise}[Attention under TP]
Extend \code{TensorParallelMLP} to a \code{TensorParallelAttention} module (heads split across
ranks, column-parallel QKV, row-parallel output). Verify against \code{MultiHeadAttention} from
\code{llmsys.layers} with gloo on two processes. What constraint on the number of heads arises?
How does GQA (\cref{ch:modern}) complicate it when $t$ exceeds the number of KV heads?
\end{exercise}

\begin{exercise}[Bubbles]
Using \code{llmsys.dist.pipeline}, plot bubble fraction and peak in-flight micro-batches for
GPipe and 1F1B with $p=8$ and $m\in\{8,16,32,64\}$. Then implement an interleaved schedule with
$v=2$ and confirm that the bubble roughly halves.
\end{exercise}

\begin{exercise}[Causal load balance in ring attention]
With $c$ devices and contiguous chunks, count the number of (query block, key block) pairs each
device computes under a causal mask. Show that the zig-zag assignment equalizes the work.
\end{exercise}

\begin{exercise}[Plan a run]
Plan the parallel layout for a 405B dense model ($L=126$, $h=16{,}384$, $s=8192$) on 16{,}384
H100s with a 16M-token global batch. Report TP, PP, CP, DP, per-GPU memory, bubble and the
exposed DP communication, stating every assumption.
\end{exercise}

% ======================================================================
\section{Lab 14: Distributed training from scratch}\label{sec:dist:lab}

\begin{labbox}{14}{Distributed training from scratch}
\labdir{lab14_distributed}

\labsection{Goals}
Implement the core distributed-training primitives yourself, verify that each is
\emph{mathematically equivalent} to single-process training, and connect measured behavior to
the cost models of this chapter.

\labsection{Setup}
No GPUs needed. \code{llmsys.dist.launch.run\_distributed(fn, world\_size)} spawns CPU processes
with the gloo backend; every collective you use also exists in NCCL. Run the tests with
\code{LAB\_IMPL=starter python -m pytest -q labs/lab08\_distributed} (one process group of 3
ranks runs all distributed checks; a non-power-of-two world catches indexing bugs).

\labsection{Tasks}
\begin{enumerate}
\item \ident{naive_allreduce}, \ident{ring_reduce_scatter}, \ident{ring_all_gather},
  \ident{ring_allreduce} using only \code{isend}/\code{irecv}; must handle any tensor shape and
  sizes not divisible by $p$.
\item \ident{allreduce_time} ($\alpha$--$\beta$, ring and naive), \ident{zero_memory_per_gpu}
  (stages 0--3), \ident{pipeline_bubble_fraction}.
\item \ident{BucketedDDP}: reverse-order buckets, post-accumulate-grad hooks, async all-reduce,
  \ident{finish_gradient_synchronization}. Weights after 3 steps must equal full-batch
  single-process AdamW, even when ranks start from different initializations.
\item \ident{ZeroAdamW} stage 1 and 2 with balanced \ident{partition_parameters}: same equivalence;
  each rank stores Adam state only for the parameters it owns.
\item \ident{CopyToTP}, \ident{ReduceFromTP}, \ident{ColumnParallelLinear},
  \ident{RowParallelLinear}: the tensor-parallel MLP must match the full MLP in outputs, input
  gradients and the corresponding weight-gradient slices.
\item \ident{gpipe_order} and \ident{one_f_one_b_order}: the provided simulator must report bubble
  $(p-1)/(m+p-1)$ and peak in-flight $m$ (GPipe) and $p-s$ (1F1B).
\end{enumerate}

\labsection{Measurements}
\begin{itemize}
\item \code{python bench.py allreduce --world 4}: time naive, ring and library all-reduce for
  1\,KB--16\,MB and fit $\alpha$, $\beta$ for each. Expect the library to win at small sizes
  (better latency) and your ring to beat the naive scheme at large sizes. On a multi-GPU node
  switch to \code{backend="nccl"} and compare the fitted bandwidth with NVLink specs. (Timings on
  a shared or oversubscribed CPU are noisy; take the minimum of several repetitions.)
\item \code{python bench.py ddp --world 2}: per-parameter all-reduce vs.\ bucketed DDP.
\item \code{python bench.py pipeline --stages 8 --mb 32}: compare with \cref{eq:dist:bubble}.
\end{itemize}

\labsection{Deliverables}
Passing tests; a plot of the fitted $\alpha$--$\beta$ curves against your measurements; a short
table comparing DDP variants; and a written 3D-parallel plan for a 70B model on 256 GPUs with
the napkin math of \cref{sec:dist:3d}.

\labsection{Stretch goals}
Implement ZeRO-3 for an \code{nn.Sequential} (gather each layer's parameters in a forward
pre-hook, free them afterwards, re-gather in backward); a hierarchical all-reduce; a real
two-stage pipeline with \code{dist.send}/\code{dist.recv} executing your 1F1B order; ring
attention on top of \code{llmsys.flash} from \cref{ch:attention}.
\end{labbox}
